% !TEX root = ../1-main.tex
\chapter{Implementation Code}
% Per-appendix pagination: this appendix's pages are C-1, C-2, ...
\setcounter{page}{1}
\label{app:code}

This appendix reproduces the code that implements the corpus-steered query
expansion (CSQE) pipeline described in Section~\ref{sec:meth_csqe}: the system
prompts used for both the blind and the corpus-grounded generation branches, the
routine that assembles the final expanded query, and the configuration that fixes
every hyperparameter reported in Chapter~\ref{chap:methodology}. The listings are
taken verbatim from the experiment notebooks, with diagnostic print statements and
environment-specific paths removed for readability. The complete code base,
including the retrieval, fusion and evaluation components, is available in the
repository given in Section~\ref{app:repo}.

%======================================================================
\section{Query Expansion System Prompts}
\label{app:code_prompts}
%======================================================================

Two system prompts are used. The first is the blind (Query2Doc-style) prompt of
Section~\ref{sec:meth_q2d_llm}, which asks the model to write a passage answering
the query from its own parametric knowledge. The second is the corpus-grounded
CSQE prompt of Section~\ref{sec:meth_csqe}, which instead casts the model as an
information-retrieval assistant and asks it to extract key sentences from the
passages returned by the first pass. Both require Arabic output.

\begin{codeblock}
\begin{lstlisting}[caption={Blind and corpus-grounded system prompts.},label={lst:systems}]
BLIND_SYSTEM = (
    "You are asked to write a passage that answers the given query. "
    "Do not ask the user for further clarification. "
    "Respond in Arabic only."
)

CSQE_SYSTEM = (
    "You are an information retrieval assistant. "
    "You will examine retrieved documents and extract key sentences "
    "relevant to the query. "
    "The documents are in Arabic. Respond in Arabic only."
)
\end{lstlisting}
\end{codeblock}

The corpus-grounded branch is additionally prefixed with a single English worked
example, taken from Table~2 of the original CSQE paper \cite{lei_2024_csqe}. This
is the sole departure from the zero-shot protocol described in
Section~\ref{sec:meth_q2d_modifications}; Aya Expanse, being multilingual,
transfers the pattern to Arabic queries while still producing Arabic output.

\begin{codeblock}
\begin{lstlisting}[caption={The one-shot demonstration prepended to the corpus-grounded prompt.},label={lst:oneshot}]
CSQE_ONE_SHOT = (
    'Query: "how are some sharks warm blooded"\n'
    'Retrieved documents:\n'
    '1. Most sharks are cold-blooded. Some, like the Mako and the Great '
    'white shark, are partially warm-blooded (they are endotherms)...\n'
    '2. Are sharks cold-blooded or warm-blooded? Sharks have a reputation '
    'as cold-blooded...\n'
    '3. Great white sharks are some of the only warm blooded sharks...\n'
    'You will begin by examining the initially retrieved documents and '
    'identifying the ones that are relevant, even partially, to the query. '
    'Once the relevant documents are identified, you will extract the key '
    'sentences from each document that contribute to their relevance.\n'
    'Based on the query "how are some sharks warm blooded", I have examined '
    'the initially retrieved documents. Here are the relevant documents and '
    'the key sentences extracted from each:\n'
    'Document 1:\n'
    '"Most sharks are cold-blooded. Some, like the Mako and the Great white '
    'shark, are partially warm-blooded (they are endotherms)."\n'
    'Document 3:\n'
    '"Great white sharks are some of the only warm-blooded sharks."\n'
)
\end{lstlisting}
\end{codeblock}

The two prompt builders below assemble the final user messages. The blind branch
passes the query through unchanged, because its system prompt already carries the
task; the corpus-grounded branch numbers the truncated first-pass passages and
appends the extraction instruction.

\begin{codeblock}
\begin{lstlisting}[caption={Prompt construction for the corpus-grounded and blind branches.},label={lst:builders}]
def build_csqe_prompt(query, retrieved_docs_truncated):
    """CSQE corpus-grounded prompt. The model should EXTRACT sentences
    from the retrieved documents, not generate freely."""
    docs_str = '\n'.join(
        f'{i+1}. {doc}' for i, doc in enumerate(retrieved_docs_truncated)
    )
    prompt = (
        f'{CSQE_ONE_SHOT}\n'
        f'Query: "{query}"\n'
        f'Retrieved documents:\n{docs_str}\n'
        f'You will begin by examining the initially retrieved documents '
        f'and identifying the ones that are relevant, even partially, to '
        f'the query. Once the relevant documents are identified, you will '
        f'extract the key sentences from each document that contribute to '
        f'their relevance. Respond in Arabic only.'
    )
    return prompt


def build_blind_prompt(query):
    """Blind (Query2Doc) prompt. The system prompt handles the task;
    the user message is just the query."""
    return query
\end{lstlisting}
\end{codeblock}

%======================================================================
\section{Query Repetition and Expansion Assembly}
\label{app:code_assembly}
%======================================================================

The routine below is the core of the method. For a batch of queries it builds both
prompt types, invokes the generator four times (twice corpus-grounded, twice
blind), and assembles the final expanded query. The assembly step is the single
line that implements the repetition factor $\alpha$ of
Section~\ref{sec:meth_repetition}: the original query is repeated $\alpha$ times
and concatenated with all four generated passages, so that the original query
terms are not diluted by the much longer expansion when the result is scored by
BM25.

\begin{codeblock}
\begin{lstlisting}[caption={CSQE batch enhancement and final query assembly.},label={lst:assembly}]
def batch_enhance(self, qids, queries):
    """Full CSQE pipeline for a batch of queries.
    Four batched generation passes: corpus x2 + blind x2."""
    n = len(qids)
    corpus_bs     = self.config['corpus_batch_size']
    blind_bs      = self.config['blind_batch_size']
    max_corpus_len = self.config['max_corpus_prompt_length']
    max_blind_len  = self.config['max_blind_prompt_length']

    # Step 1: build all prompts
    corpus_prompts, blind_prompts, all_retrieved = [], [], []
    for qid, query in zip(qids, queries):
        docs = self.get_retrieved_docs(qid)      # BM25 first pass, truncated
        all_retrieved.append(docs)
        doc_texts = [d['text'] for d in docs if d['text']]
        corpus_prompts.append(build_csqe_prompt(query, doc_texts))
        blind_prompts.append(build_blind_prompt(query))

    # Step 2: corpus-grounded expansions (2 samples)
    corpus_samples_1 = self.batch_generate(
        CSQE_SYSTEM, corpus_prompts, corpus_bs,
        temperature=1.0, max_length=max_corpus_len)
    corpus_samples_2 = self.batch_generate(
        CSQE_SYSTEM, corpus_prompts, corpus_bs,
        temperature=1.0, max_length=max_corpus_len)

    # Step 3: blind expansions (2 samples)
    blind_samples_1 = self.batch_generate(
        BLIND_SYSTEM, blind_prompts, blind_bs,
        temperature=1.0, max_length=max_blind_len)
    blind_samples_2 = self.batch_generate(
        BLIND_SYSTEM, blind_prompts, blind_bs,
        temperature=1.0, max_length=max_blind_len)

    # Step 4: assemble the expanded query
    alpha = self.config['query_repetition']
    results = []
    for i in range(n):
        corpus_exps = [corpus_samples_1[i], corpus_samples_2[i]]
        blind_exps  = [blind_samples_1[i],  blind_samples_2[i]]
        all_exps    = corpus_exps + blind_exps

        # Repeat the original query alpha times, then append all
        # four generated passages. This is the expanded query that
        # is passed to the BM25 second pass.
        final = (queries[i] + ' ') * alpha + ' '.join(all_exps)

        results.append({
            'qid':                qids[i],
            'original':           queries[i],
            'retrieved_docids':   [d['docid'] for d in all_retrieved[i]],
            'corpus_expansions':  corpus_exps,
            'blind_expansions':   blind_exps,
            'enhanced':           final,
        })
    return results
\end{lstlisting}
\end{codeblock}

%======================================================================
\section{CSQE Configuration}
\label{app:code_config}
%======================================================================

The configuration below fixes every hyperparameter of the primary CSQE experiment
reported in Section~\ref{sec:res_csqe}. The values correspond to those stated in
Section~\ref{sec:meth_csqe}: a first pass of $k=5$ documents each truncated to 128
tokens, two corpus-grounded and two blind samples, a repetition factor of
$\alpha = 4$, and a generation temperature of 1.0 chosen to give diversity across
the four samples.

\begin{codeblock}
\begin{lstlisting}[caption={Configuration for the primary CSQE experiment.},label={lst:config}]
CONFIG = {
    # BM25 first pass
    'top_k_docs':            5,    # top-5 carry most of the signal and
                                   # halve the corpus prompt length
    'doc_truncation_tokens': 128,  # truncate each retrieved doc

    # LLM generation
    'model_name':          'CohereForAI/aya-expanse-8b',
    'temperature':         1.0,    # diversity across multiple samples
    'max_new_tokens':      128,
    'top_p':               0.9,
    'num_corpus_samples':  2,      # corpus-grounded expansions
    'num_blind_samples':   2,      # blind (Query2Doc-style) expansions

    # Query construction: alpha = num_corpus + num_blind = 4
    'query_repetition':    4,

    # Batching (cross-query batching on the A100)
    'corpus_batch_size':        8,     # long prompts (~800-1000 tokens)
    'blind_batch_size':         32,    # short prompts (~60-80 tokens)
    'max_corpus_prompt_length': 2048,
    'max_blind_prompt_length':  512,
    'macro_batch_size':         200,   # queries per checkpoint cycle
}
\end{lstlisting}
\end{codeblock}

%======================================================================
\section{Source Code Repository}
\label{app:repo}
%======================================================================

The complete implementation is publicly available at:

\begin{center}
\url{https://github.com/Osmanoor/arabic-rag-query-enhancement}
\end{center}

\noindent The repository contains the components not reproduced above, including
the BM25S and mDPR retrieval wrappers, the reciprocal rank fusion and convex
combination implementations used for the hybrid experiments of
Section~\ref{sec:res_hybrid}, the \texttt{pytrec\_eval} evaluation harness that
produced every metric reported in Chapter~\ref{chap:results}, the per-model query
generation notebooks, and the ablation and error-analysis notebooks.
